\section{Risk-Controlled Safe Stopping via Policy-Level Certification}
\label{sec:rcs}

We first define the risk conditional on stopping and show how to test each candidate policy using the calibration samples on which it stops. 
We then combine these tests to certify a set of policies and select one based on a user-chosen utility, defaulting to full acquisition if none can be certified. 
Finally, we extend the procedure to class-conditional false negative and false positive rates and describe the stopping rules used in our experiments.

% \gls{method} treats each candidate stopping rule as a complete trajectory-level policy, measures the loss it incurs on the calibration trajectories it stops on, and uses the LTT machinery of \cref{sec:ltt} to certify which rules control the conditional risk.

\subsection{Conditional risk of a stopping policy}
Consider a family of stopping rules $\cD_\textsc{stop}(\mbx_t; \lambda) \in \{0, 1\}$ indexed by a finite grid $\lambda \in \Lambda = \{\lambda_1, \dots, \lambda_m\}$. Each rule induces a stopping time, the first step at which it fires,
\begin{align}
\label{eq:stopping_time}
    \tau_\lambda(\mbx) = \inf\{t \in \{1, \dots, d \} \mid \cD_\textsc{stop}(\mbx_t; \lambda) = 1\},
\end{align}
with $\tau_\lambda(\mbx) = \infty$ if the rule never fires, in which case the policy defers to full acquisition. When it fires at $t \le d$, the system commits to $\hat{y}(\mbx_t)$. For a bounded loss $\ell(y, \hat{y}) \in [0, B]$, the quantity we control is the expected loss given that the policy stops,
\begin{align}
\label{eq:cond_risk_def}
    R_{\text{cond}}(\lambda) = \E_{\mbx, y \sim \cP^{\text{data}}}\left[\ell\left(y, \hat{y}(\mbx_{\tau_\lambda(\mbx)})\right) \;\middle|\; \tau_\lambda(\mbx) < \infty \right].
\end{align}
Our goal is to select a $\hat{\lambda} \in \Lambda$ with $R_{\text{cond}}(\hat\lambda) \le \alpha$ with probability at least $1-\delta$ over the draw of $\cD_{\text{cal}}$. 
Changing $\lambda$ changes not only \emph{which} samples the policy stops on, but also changes \emph{when} it stops. 
It therefore also changes the partial observation each prediction is made. 
The map $\lambda \mapsto R_{\text{cond}}(\lambda)$ can therefore jump and need not be monotone.

\paragraph{Testing each policy on the samples it stops.}
For each $\lambda$, let $\cS(\lambda) = \{ i \le n \mid \tau_\lambda(\mbx_i) < \infty \}$ be the indices of the calibration samples where the policy stops, and let $n(\lambda) = |\cS(\lambda)|$. 
The empirical conditional risk is the mean loss over this subset,
\begin{align}
    \widehat{R}_{\text{cond}}(\lambda) = \frac{1}{n(\lambda)} \sum_{i \in \cS(\lambda)} \ell\left(y^i, \hat{y}(\mbx^i_{ \tau_\lambda(\mbx^i)})\right),
\end{align}
and we test $H_{0,\lambda} : R_{\text{cond}}(\lambda) > \alpha$.
The set $\cS(\lambda)$ is selected by the very rule under test. 
Because the calibration trajectories are i.i.d.\ and membership in $\cS(\lambda)$ depends only on the trajectory itself, the stopped samples are conditional on $n(\lambda)$ draws from the distribution conditioned on stopping. 
Hoeffding's inequality then yields a p-value that is valid conditional on $n(\lambda)$,
\begin{align}
\label{eq:hoeffding_pvalue}
    p(\lambda) =
    \begin{cases}
        1, & \text{if } n(\lambda) = 0 \text{ or } \widehat{R}_{\text{cond}}(\lambda) \ge \alpha, \\[2pt]
        \exp\!\left(-2 n(\lambda)\big(\alpha - \widehat{R}_{\text{cond}}(\lambda)\big)^2 / B^2\right), & \text{otherwise}.
    \end{cases}
\end{align}
This conditional validity also makes the p-value valid over the random draw of the calibration set. 
%and validity conditional on $n(\lambda)$ implies validity marginally over the calibration set.

\paragraph{Multiplicity correction and selection.}
Testing $m$ hypotheses at once calls for family-wise error control. We apply Holm's step-down procedure \citep{holm1979simple}: sort the p-values $p_{(1)} \le \cdots \le p_{(m)}$, reject $H_{(j)}$ as long as $p_{(j)} \le \delta/(m-j+1)$, and stop at the first failure. The rejected hypotheses form the safe set $\widehat{\Lambda}_{\text{safe}}$, and every $\lambda$ in it controls the conditional risk. We deploy the safe rule that maximizes a user-chosen utility $U(\lambda)$, typically the fraction of samples stopped or the negative mean stopping time. If $\widehat{\Lambda}_{\text{safe}}$ is empty, \gls{method} reports that no candidate could be certified at the chosen confidence level using the available calibration data, and every sample proceeds to full acquisition. This does not establish that no safe candidate exists or certify the risk of the full-acquisition predictor.

\begin{corollary}[Conditional risk control; corollary of LTT]
    \label{thm:conditional_risk}
    Let $\cD_{\text{cal}}$ consist of $n$ i.i.d.\ sample trajectories and let $\hat{\lambda}$ be selected by the procedure above. Then, provided $\widehat{\Lambda}_{\text{safe}} \neq \emptyset$,
    \begin{align}
        \bbP_{\cD_{\text{cal}}}\left[ \E_{\mbx, y} \left[ \ell\left(y, \hat{y}(\mbx_{\tau_{\hat\lambda}(\mbx)})\right) \;\middle|\; \tau_{\hat{\lambda}}(\mbx) < \infty \right] \leq \alpha \right] \geq 1 - \delta.
    \end{align}
\end{corollary}
The result is a direct application of LTT \citep{angelopoulos2025learn}: each $\lambda$ is a selective predictor on trajectories, and LTT already certifies a selective risk over a data-dependent selected subset. 
The argument requires a policy family and selection events fixed independently of the risk-calibration labels, and independent calibration trajectories drawn from the same induced distribution as deployment trajectories. The guarantee concerns the eventually stopped cohort $\{\tau_{\hat\lambda}<\infty\}$; it does not assert simultaneous risk control for the cumulative cohorts $\{\tau_{\hat\lambda}\le t\}$ at every checkpoint.
%The argument uses nothing about the selection event beyond its being fixed before calibration, so it applies verbatim to any pre-specified event. 
Pseudocode is given in \cref{eq:alg_ltt} in the appendix.

\subsection{Class-conditional FNR and FPR}
\label{sec:fnr_fpr}
\Cref{eq:cond_risk_def} covers mis-classification directly. Clinical practice asks instead for class-conditional rates: the false negative rate among truly positive cases and the false positive rate among truly negative ones, rather than a pooled or prevalence-weighted error. For a task with labels $y \in \{0,1\}^K$ (binary when $K=1$) we certify, for each label $k$,
\begin{align}
    \mathrm{FNR}_k(\lambda) = \bbP\!\left(\hat{y}_k(\mbx_{\tau_\lambda}) = 0 \;\middle|\; y_k = 1,\ \tau_\lambda(\mbx) < \infty\right) \le \alpha_k,
\end{align}
and $\mathrm{FPR}_k$ symmetrically. This is the conditional risk of \cref{eq:cond_risk_def} under the refined selection event $\{y_k = 1,\ \tau_\lambda < \infty\}$: the effective calibration set is the stopped positives $\cS^+_k(\lambda) = \{ i : \tau_\lambda(\mbx^i) < \infty,\ y^{i}_k=1 \}$ and the loss is $\mathbf{1}\{\hat y^i_{k}=0\}$. Because the loss is binary we replace Hoeffding with the exact binomial tail, which is tighter, and we combine labels by $p(\lambda) = \max_k p_k(\lambda)$, which is valid under arbitrary dependence between labels and accommodates label-specific targets $\alpha_k$. The policy family, the stopping rule, and the Holm correction are unchanged. Class-conditional certification requires sufficient stopped positive (or negative) calibration examples; a policy with none cannot be certified for that class.

\paragraph{Policy families.}
\gls{method} places no restriction on the form of the stopping rule beyond the family being fixed before calibration. We use three. A \emph{confidence threshold} on the softmax score (\textsc{cifar}-10); the \emph{\textsc{jafa} Q-margin}, the gap between the Q-value of ``stop'' and that of the best acquisition action (\cref{sec:experiments_2}); and \emph{ambiguity bands} for the \textsc{mri} tasks, where $\lambda \in [0,1]$ scales a per-label uncertainty band from never stopping ($\lambda = 0$) to an aggressive data-driven band ($\lambda = 1$), and the rule fires once every label's score has left its band (\cref{app:ambiguity_bands}).
